The results of Chapter 6 are obtained under a \emph{within-subject} protocol: for
each participant, a model is trained on 280 of that participant's trials and tested
on the remaining 120. Limitation~2 of this thesis acknowledges that such figures
cannot support a generalisation claim. This chapter removes that limitation by
re-running the entire pipeline under a subject-disjoint protocol and reporting what
the system achieves on participants it has never seen.

\section{Motivation: What Within-Subject Evaluation Conceals}
\label{sec:cs_motivation}

Under the within-subject protocol, every encoder sees the test participant's face,
voice and neural signature during training. The reported accuracy therefore answers
the question \emph{``how well does the system recognise emotions in a person it has
already studied?''}, not \emph{``does the system recognise emotions?''}. For a
deployment scenario---clinical screening, human-computer interaction, or the
demonstration component of this thesis---the second question is the operative one,
because every new user is by definition a participant the system has never trained on.

This distinction is not merely theoretical. A facial-emotion encoder trained and
tested on the same individual can, in principle, minimise its training loss by
learning \emph{this particular person's} angry face rather than a generalisable
representation of anger. Such a model would report high within-subject accuracy while
carrying no transferable knowledge. Distinguishing the two cases requires a protocol
in which the test participant is absent from training entirely.

Three trimodal-fusion methods have been published on
EAV---AMERL~\cite{Yin2024AMERL}, Hyper-MML~\cite{Kang2025HyperMML}, and
EEG-MoCE~\cite{Anon2026EEGMoCE}---and all three report within-subject results.
To our knowledge, no subject-independent evaluation of any trimodal fusion method on
the EAV benchmark has been published. The evaluation reported in this chapter is
therefore the first of its kind on this corpus.

\section{Protocol: Five-Fold Subject-Wise Cross-Validation}
\label{sec:cs_protocol}

\subsection{Fold Construction}

The 42 participants are partitioned into five disjoint groups by a fixed random seed,
producing folds of 9, 9, 8, 8 and 8 participants. For fold $k$, the participants in
group $k$ are held out entirely and the remaining participants constitute the
training population. Because each EAV participant contributes a class-balanced 400
trials (80 per emotion), no stratification is required at the participant level: any
partition of participants is automatically class-balanced.

The partition is deterministic given the seed, so the same assignment is reproduced on
every machine without coordination---a practical requirement, since the five folds
were executed in parallel across three separate GPU sessions. Table~\ref{tab:cs_folds}
reports the assignment.

\begin{table}[htbp]
\centering
\caption{Subject-wise fold assignment. Each participant is held out exactly once.}
\label{tab:cs_folds}
\begin{tabular}{clc}
\hline
\textbf{Fold} & \textbf{Held-out participants} & \textbf{Training participants} \\
\hline
0 & 5, 8, 19, 22, 26, 27, 30, 40, 42 & 33 \\
1 & 6, 17, 21, 25, 29, 31, 33, 38, 39 & 33 \\
2 & 7, 11, 12, 16, 18, 24, 32, 36 & 34 \\
3 & 1, 4, 10, 13, 20, 23, 28, 35 & 34 \\
4 & 2, 3, 9, 14, 15, 34, 37, 41 & 34 \\
\hline
\end{tabular}
\end{table}

Five folds are used rather than leave-one-subject-out (LOSO). LOSO would require 42
folds $\times$ 3 encoders $=$ 126 encoder trainings; at the measured cost of
approximately 2.7 hours per fold, this is approximately 23 hours of GPU time for
five folds against roughly 190 hours for LOSO. Five-fold subject-wise
cross-validation retains the essential property---every participant is tested exactly
once as an unseen individual---at a fraction of the cost. Because every participant is
held out exactly once, a complete run yields 42 per-participant accuracies per method,
the same shape as the within-subject results table, permitting a paired comparison
between the two protocols.

\subsection{The Encoder-Reuse Constraint}
\label{sec:cs_leakage}

The single most important methodological constraint in this chapter is that the
cached artefacts from the within-subject rollout \emph{cannot} be reused.

Each per-subject encoder state dictionary was fine-tuned on that subject's own
training trials, and each cached feature matrix was produced by that encoder.
Training a fusion head on participants 1--41 and testing on participant 42 using
those caches would evaluate on features generated by a model that had already been
fitted to participant 42's data. The resulting figure would be labelled
``cross-subject'' while containing exactly the contamination the protocol exists to
eliminate, and---critically---it would look entirely plausible, since it would fall
somewhere between the within-subject and true cross-subject values.

Consequently, all three per-modality encoders are retrained from scratch for every
fold, using only that fold's training participants. This retraining is the entire
computational cost of the cross-subject evaluation and is not optional.

\subsection{Inner Validation and Model Selection}
\label{sec:cs_innerval}

A second, subtler leakage channel arises from model selection. The per-modality
trainers perform best-epoch checkpointing, and the fusion heads require an epoch at
which to stop. Using the held-out fold for either purpose would constitute selection
on the test set.

Three participants are therefore withheld from each fold's training population and
reserved exclusively for model selection, leaving 30 or 31 participants for parameter
fitting. The split is made at the \emph{participant} level so that inner validation
mirrors the outer protocol: the selection signal comes from individuals the encoders
have never seen, exactly as the test fold does.

The same three participants later serve as the validation set for the fusion heads.
This is deliberate. Their cached features were produced by encoders that never saw
them, which is precisely the condition under which the test fold's features are
generated. Validating the fusion head on the fitting participants instead would be
optimistically biased, because their features come from encoders trained on them.

\subsection{Trial Allocation and Feature Standardisation}

Under a subject-wise protocol the within-subject 280/120 split has no meaning: a
training participant contributes every trial they have, and a held-out participant is
scored on every trial they have. Each fold therefore trains on approximately
13,200--13,600 trials---roughly 47 times the 280 trials available to a per-subject
model---and each held-out participant is evaluated on all 400 of their trials.

Per-modality features are standardised per fold, with means and variances computed on
the training participants only. Cross-subject feature distributions drift in a way
that within-subject distributions do not, and the three modalities have different
dimensionalities (768, 768 and 960) and different natural scales; without
standardisation, a modality could dominate the fusion input purely through the
magnitude of its activations.

\section{Implementation}
\label{sec:cs_implementation}

The cross-subject pipeline mirrors the three-stage structure of
Section~\ref{sec:cs_protocol} and is implemented in \texttt{cv\_pipeline.py}, with the
fold partition isolated in \texttt{cv\_split.py} so that it can be independently
validated and shared across parallel sessions.

\textbf{Stage A---Fold-wise encoder training.} For each fold, AST, ViT and EEGNet are
trained on that fold's fitting participants. Per-modality epoch budgets follow the
within-subject configuration (AST 2+2, ViT 3+2, EEGNet up to 150 with early stopping),
reduced for AST because each fold presents approximately 47 times more data per epoch.

\textbf{Stage B---Feature and logit extraction.} Each fold's encoders extract
pre-classifier features \emph{and} per-modality class logits for all 42 participants.
Extracting logits alongside features in the same forward pass avoids the post-hoc
regeneration step that was required during the within-subject rollout
(Section~\ref{sec:c9}). Stage B costs approximately 54 minutes per fold.

\textbf{Stage C---Fusion training and scoring.} All fusion variants are trained on the
fold's fitting participants and evaluated per held-out participant. Because Stage C
reads only cached features, the entire fusion analysis can be re-run across all five
folds in minutes without touching a GPU.

Eight methods are scored: the three per-modality encoders evaluated independently
(\texttt{audio\_only}, \texttt{vision\_only}, \texttt{eeg\_only}), naive late fusion,
cross-attention fusion, concat-MLP fusion, modality-dropout fusion at full modality,
and the modality-dropout model evaluated with EEG zeroed.

\section{Results}
\label{sec:cs_results}

Table~\ref{tab:cs_main} reports the 42-participant cross-subject results alongside
the within-subject figures of Chapter~6.

\begin{table}[htbp]
\centering
\caption{Cross-subject results across 42 participants (5-fold subject-wise
cross-validation), compared with the within-subject protocol. Chance is 20.0\%.}
\label{tab:cs_main}
\begin{tabular}{lccccc}
\hline
\textbf{Method} & \textbf{Cross-subj.} & \textbf{SD} & \textbf{Min--Max} & \textbf{Within-subj.} & \textbf{$\Delta$} \\
\hline
Naive late fusion      & \textbf{62.15\%} & 0.093 & 38.3--78.2 & 77.5\% & $-15.3$ pp \\
Concat-MLP fusion      & 58.80\% & 0.089 & 39.8--77.5 & 81.7\% & $-22.9$ pp \\
Modality dropout (full)& 58.77\% & 0.110 & 29.7--77.7 & 84.7\% & $-25.9$ pp \\
Modality dropout (AV)  & 57.89\% & 0.110 & 28.5--76.5 & 81.5\% & $-23.6$ pp \\
Cross-attention fusion & 57.61\% & 0.111 & 29.5--81.0 & 80.2\% & $-22.6$ pp \\
Audio only (AST)       & 57.51\% & 0.086 & 40.7--73.0 & 57.1\% & $\mathbf{+0.4}$ \textbf{pp} \\
Vision only (ViT)      & 41.18\% & 0.094 & 22.0--58.3 & 74.6\% & $\mathbf{-33.4}$ \textbf{pp} \\
EEG only (EEGNet)      & 38.17\% & 0.096 & 22.8--57.0 & 43.9\% & $-5.7$ pp \\
\hline
\end{tabular}
\end{table}

\subsection{Fold Consistency}

Before interpreting the aggregate figures, it is necessary to confirm that no single
fold dominates them. Table~\ref{tab:cs_perfold} reports per-fold means.

\begin{table}[htbp]
\centering
\caption{Per-fold mean accuracy. The narrow spread across folds indicates that the
aggregate results are not driven by a single partition.}
\label{tab:cs_perfold}
\begin{tabular}{lccccc}
\hline
\textbf{Method} & \textbf{Fold 0} & \textbf{Fold 1} & \textbf{Fold 2} & \textbf{Fold 3} & \textbf{Fold 4} \\
\hline
Naive late fusion       & 0.629 & 0.616 & 0.655 & 0.586 & 0.621 \\
Concat-MLP              & 0.589 & 0.600 & 0.601 & 0.567 & 0.582 \\
Modality dropout (full) & 0.603 & 0.610 & 0.604 & 0.565 & 0.552 \\
Modality dropout (AV)   & 0.578 & 0.598 & 0.594 & 0.552 & 0.569 \\
Cross-attention         & 0.567 & 0.597 & 0.629 & 0.564 & 0.522 \\
Audio only              & 0.570 & 0.567 & 0.599 & 0.541 & 0.600 \\
Vision only             & 0.419 & 0.396 & 0.444 & 0.407 & 0.393 \\
EEG only                & 0.345 & 0.396 & 0.400 & 0.405 & 0.366 \\
\hline
\end{tabular}
\end{table}

Naive late fusion ranges from 0.586 to 0.655 across folds and audio from 0.541 to
0.600. No fold is an outlier on any method, and the ordering of methods is stable
across folds. The protocol behaves consistently.

\subsection{Statistical Comparison of Fusion Methods}

Paired Wilcoxon signed-rank tests across the 42 per-participant accuracies are
reported in Table~\ref{tab:cs_wilcoxon}.

\begin{table}[htbp]
\centering
\caption{Paired Wilcoxon signed-rank tests on cross-subject per-participant
accuracies ($n = 42$, two-sided).}
\label{tab:cs_wilcoxon}
\begin{tabular}{llcc}
\hline
\textbf{Method A} & \textbf{Method B} & \textbf{$\bar{\Delta}$} & \textbf{$p$} \\
\hline
Naive late fusion & Cross-attention        & $+4.54$ pp & $0.0003$ \\
Naive late fusion & Modality dropout (full)& $+3.38$ pp & $0.0026$ \\
Naive late fusion & Concat-MLP             & $+3.35$ pp & $0.0015$ \\
Naive late fusion & Audio only             & $+4.64$ pp & $0.0003$ \\
\hline
Modality dropout (full) & Concat-MLP       & $-0.03$ pp & $0.717$ \\
Modality dropout (full) & Cross-attention   & $+1.16$ pp & $0.141$ \\
Concat-MLP & Cross-attention               & $+1.19$ pp & $0.204$ \\
Modality dropout (full) & Modality dropout (AV) & $+0.89$ pp & $0.090$ \\
Modality dropout (full) & Audio only       & $+1.26$ pp & $0.250$ \\
\hline
Audio only & Vision only                   & $+16.33$ pp & $<0.0001$ \\
Vision only & EEG only                     & $+3.01$ pp & $0.086$ \\
\hline
\end{tabular}
\end{table}

Two results stand out. First, \textbf{naive late fusion significantly outperforms
every learned fusion method} under the subject-independent protocol, by margins of
3.35 to 4.54 percentage points, all at $p < 0.003$. This reverses the within-subject
ordering, in which every learned method beat the parameter-free baseline.

Second, \textbf{the three learned fusion methods are statistically indistinguishable
from one another} (all pairwise $p > 0.14$). Cross-attention, concat-MLP and
modality-dropout training---which differ substantially in parameter count, inductive
bias and training discipline---converge to the same cross-subject performance. The
choice of learned fusion operator, which Chapter~6 found to be the smallest of the
within-subject design decisions, becomes irrelevant entirely under subject shift.

\subsection{Fusion Gain Over the Best Single Modality}

A stricter test than comparing fusion methods to each other is to compare each method
against the best single modality \emph{for that participant}, computed per
participant rather than on population means. Table~\ref{tab:cs_gain} reports this
comparison.

\begin{table}[htbp]
\centering
\caption{Fusion gain over the per-participant best unimodal classifier
($n = 42$, paired Wilcoxon).}
\label{tab:cs_gain}
\begin{tabular}{lccc}
\hline
\textbf{Method} & \textbf{Mean gain} & \textbf{$p$} & \textbf{Participants improved} \\
\hline
Naive late fusion       & $+3.70$ pp & $0.0020$ & 31 / 42 \\
Concat-MLP              & $+0.35$ pp & $0.488$   & 23 / 42 \\
Modality dropout (full) & $+0.32$ pp & $0.528$   & 24 / 42 \\
Cross-attention         & $-0.85$ pp & $0.887$   & 23 / 42 \\
\hline
\end{tabular}
\end{table}

Only naive late fusion delivers a statistically reliable gain over the best single
modality. None of the three learned methods does; cross-attention is marginally
\emph{below} the per-participant best unimodal classifier. Under subject-independent
evaluation, the learned fusion heads fail to justify their existence on accuracy
grounds.

\subsection{Per-Modality Transfer}
\label{sec:cs_permodality}

The per-modality rows of Table~\ref{tab:cs_main} contain the chapter's most
substantive finding. The loss incurred by moving to unseen participants is not
distributed evenly across modalities:

\begin{itemize}
    \item \textbf{Audio transfers at parity.} 57.51\% cross-subject against 57.1\%
    within-subject---a difference of $+0.4$ percentage points, well inside the noise
    of the measurement. The AST encoder's representation of emotional speech carries
    over to new speakers essentially unchanged.
    \item \textbf{EEG loses moderately.} 38.17\% against 43.9\%, a drop of 5.7
    percentage points, while remaining well above the 20\% chance level and above the
    EAV paper's published within-subject EEGNet baseline of 36.7\%.
    \item \textbf{Vision collapses.} 41.18\% against 74.6\%, a drop of 33.4
    percentage points. The strongest modality under within-subject evaluation becomes
    the second weakest under subject-independent evaluation.
\end{itemize}

The ordering of modalities inverts. Within-subject, vision (74.6\%) $>$ audio (57.1\%)
$>$ EEG (43.9\%); cross-subject, audio (57.5\%) $>$ vision (41.2\%) $>$ EEG (38.2\%).
Audio, not vision, is the strongest single modality on unseen participants, and the
margin is large and highly significant ($+16.33$ pp, $p < 0.0001$).

A supplementary observation reinforces the interpretation. During calibration, the
vision encoder was trained under two budgets: the P2-matched 3~frozen $+$
2~fine-tuning epochs, and a reduced 1$+$1 schedule. Training the vision encoder
\emph{longer} produced \emph{worse} cross-subject accuracy (41.9\% against 45.8\% on
fold 0, $p = 0.031$), while the same additional training is beneficial within-subject.
Additional gradient steps on the training participants' faces actively degrade
performance on new faces.

\section{Discussion}
\label{sec:cs_discussion}

\subsection{Identity Inflation in Facial-Emotion Recognition on EAV}

The asymmetry between audio ($+0.4$ pp) and vision ($-33.4$ pp) admits a direct
interpretation: a substantial portion of EAV's within-subject vision accuracy
reflects participant identity rather than emotional expression.

Facial appearance is strongly identity-bound. When training and test trials come from
the same individual, a ViT can reduce its loss by associating specific facial
configurations---this person's particular smile, brow geometry, or resting
expression---with emotion labels. Such a model performs excellently on that individual
and poorly on anyone else. Vocal prosody carries far less identity-specific structure
relevant to the emotion label: raised pitch and increased energy signal arousal in a
broadly speaker-invariant manner, which is consistent with audio's complete transfer.

The finding that additional vision training degrades cross-subject accuracy supports
this reading. If the drop reflected underfitting, more training would help; instead it
hurts, which is the expected signature of an encoder fitting participant-specific
structure.

This reframes the 74.6\% vision figure reported in Chapter~6. That number is correct
as measured, but it should be understood as the performance of a
participant-specialised model, not as evidence that facial-emotion features on EAV
generalise. The implication extends beyond this thesis: any within-subject vision
result on EAV, including those of prior published work on this benchmark, is subject
to the same inflation.

\subsection{Why Learned Fusion Loses Under Subject Shift}

The reversal of the fusion ordering has a coherent mechanism. A learned fusion head
does two things: it extracts cross-modal interactions, and it learns how much to
trust each modality. The second capability is fitted to the training participants'
modality reliability---in this corpus, a population in which vision is the strongest
modality, because within the training population vision \emph{is} strong.

On a held-out participant, that fitted reliability estimate is partly wrong. Vision
does not transfer, so a fusion head that has learned to weight vision heavily inherits
a systematically miscalibrated prior. Naive late fusion, weighting all three modalities
equally by construction, has no such fitted quantity and therefore nothing that can be
wrong in this way. Parameter-free fusion is more robust precisely because it is less
expressive.

This is a protocol-dependent conclusion rather than a claim that attention-based
fusion is without merit. Under within-subject evaluation, learned fusion wins
significantly; under subject-independent evaluation, it does not. Both statements are
true, and reporting only the first would misrepresent the system's behaviour on new
users.

\subsection{The Magnitude of the Overall Drop}

The headline system accuracy falls from 84.7\% (modality dropout, within-subject) to
62.2\% (naive late fusion, cross-subject). This 22.5-percentage-point gap is the
quantity that within-subject evaluation concealed.

The gap should not be read as a defect of the method. It is a property of the
evaluation protocol, and an equivalent gap is expected to exist---though it has not
been measured---for the prior published methods on this benchmark, all of which report
within-subject figures. What the present evaluation contributes is the measurement
itself, together with a localisation of where the loss originates.

It is also worth noting what did \emph{not} degrade. The per-participant standard
deviation is 0.093 cross-subject against 0.083 within-subject: the spread of
individual difficulty is essentially unchanged. The level moves; the shape of the
population distribution does not.

\subsection{Robustness Under Missing EEG}

The zero-EEG configuration retains its practical viability. Modality-dropout fusion
scores 58.77\% at full modality and 57.89\% with EEG zeroed---a difference of 0.89
percentage points that is not statistically significant ($p = 0.090$). Within-subject,
the same comparison cost 3.2 percentage points.

The demonstration path, which operates without an EEG headset, therefore degrades
gracefully on unseen participants as well as on seen ones. The corollary is less
favourable to the trimodal premise: if removing EEG entirely costs under one
percentage point, EEG's marginal contribution to the fused prediction on unseen
participants is small, notwithstanding its 38.2\% standalone accuracy.

\section{Answers to Research Questions}
\label{sec:cs_rq}

\textbf{RQ6---Cross-subject generalisation:} The system generalises to unseen
participants at 62.2\% mean accuracy (naive late fusion), substantially above the
20\% chance level but 15.3 percentage points below its within-subject figure and 22.5
points below the best within-subject configuration. All per-modality and fusion
methods remain above chance on every fold, confirming that the pipeline learns
genuinely subject-transferable structure; the magnitude of the drop quantifies how
much of the within-subject result was participant-specific.

\textbf{RQ7---Protocol dependence of the architectural contribution:} No. Under
subject-independent evaluation, learned fusion does not outperform naive
softmax averaging; the ordering established in Chapter~6 reverses, and all three
learned fusion variants become statistically indistinguishable from one another. The
architectural advantage demonstrated within-subject is a property of that protocol and
does not survive subject shift.
